% Skabelon til PDF-udgaven af en DMA-eksamen.
% Kopiér filen (fx til reexam-april-2026.tex), og ret indstillingerne nedenfor.
% Fjern %-tegnet foran \answerstrue for at lave udgaven med svar og forklaringer.
\documentclass[12pt]{article}
\newif\ifanswers

% ===== Indstillinger for den konkrete eksamen =====
%\answerstrue % med svar og forklaringer
\newcommand{\examdate}{17. april 2026}
\newcommand{\examduration}{2 timer}
\newcommand{\problemfile}{dma-problems-april-2026} % opgavefil i denne mappe, uden .tex
% ==================================================

\usepackage{amssymb}
\usepackage{amsmath}
\usepackage{verbatim}
\usepackage{graphicx}
\usepackage{enumitem}
\usepackage{xspace}
\usepackage{tikz}
\usetikzlibrary{arrows,arrows.meta,automata,backgrounds,calc,
  decorations.pathmorphing,decorations.pathreplacing,fit,matrix,patterns,
  positioning,shapes,shapes.geometric,trees}
\usepackage{calc}
\usepackage{fancybox}
\usepackage{clrscode3e}
\usepackage{totcount}
\usepackage{geometry}
\geometry{a4paper, margin=1in}

\pagenumbering{gobble}

% Antal opgaver: hvert \spoergsmaal er en \section.
\regtotcounter{section}

\newcounter{svarside}
\setcounter{svarside}{1}
\regtotcounter{svarside}

\newcommand{\svarside}{\addtocounter{svarside}{1}%
Skriv dit svar på side \thesvarside\ i Word dokumentet.
}

\newcommand{\punkt}[2]{\vspace{.3cm}\noindent\begin{Sbox}
    \begin{minipage}{16cm}\textbf{#1} \xspace {#2}\xspace\svarside\end{minipage}\end{Sbox}\fbox{\TheSbox}\\
    \vspace{.1cm}
    \noindent
    }

\tikzset{
  treenode/.style = {align=center, inner sep=0pt, text centered,
    font=\sffamily,text=white},
  black/.style = {treenode, circle, white, font=\sffamily\bfseries, draw=black,
    fill=black, text width=1.5em},
  red/.style = {treenode, circle, white, draw=red,
    fill=red, text width=1.5em, very thick},
  white/.style = {treenode, rectangle, draw=none,
    minimum width=0.5em, minimum height=0.5em},
  gray/.style = {treenode, rectangle, draw=black,
    minimum width=0.5em, minimum height=0.5em}
}

\newcommand{\spoergsmaal}[1]{\section{#1\mcprocent}}
\newcommand{\nylinie}{\\}
\newcommand{\multiplechoicestart}{}
\newcommand{\multiplechoiceslut}{}
\newcommand{\mulighederstart}{\begin{enumerate}}
\newcommand{\mulighederslut}{\end{enumerate}}
\newcommand{\svar}{\item}
\newcommand{\itxsvar}[1]{} % For ITX, ignore in LaTeX
\newcommand{\qtipoints}[1]{} % For QTI/WISEflow, ignore in LaTeX
\newcommand{\etsvar}{Der er præcis ét korrekt svar, men det er tilladt at vælge flere svar hvis du er i tvivl.\xspace}
\newcommand{\fleresvar}{Vælg ét eller flere korrekte svar.\xspace}
\newcommand{\mcprocent}{\xspace (4\%)}

\ifanswers
\newcommand{\korrekt}{\checkmark}
\newcommand{\medsvar}{{\bf Opgaver med vejledende svar}\\}
\newcommand{\forklaring}[1]{{\bf Svar: }{#1}}
\newcommand{\sideskift}{}
\else
\newcommand{\korrekt}{}
\newcommand{\medsvar}{}
\newcommand{\forklaring}[1]{}
\newcommand{\sideskift}{\newpage}
\fi

\begin{document}

\title{\Large \medsvar Eksamen i Diskret Matematik og Algoritmer\\
{\small NDAB15008U}}

\author{Københavns Universitet}
\date{\examdate}

\maketitle

\section*{Instruktioner}

Du har \examduration\ til at besvare \total{section} opgaver, beskrevet nedenfor.
Bedømmelsen baseres på det samlede antal point.
Hver opgave kan give op til 4 point, så en perfekt besvarelse giver \the\numexpr 4*\totvalue{section}\relax\ point.
%Opgavesættet har \pageref{page:last} nummererede sider.
% CLRS refererer til kursets lærebog, Introduction to Algorithms, 4th edition.

\medskip

Skriftlige hjælpemidler samt brug af lokalt installeret software er tilladt. Det er \emph{ikke} tilladt at tilgå internet eller at kommunikere med andre.

\medskip

Alle opgaver besvares i ITX systemets multiple-choice-modul.
I hvert multiple-choice spørgsmål er det enten angivet at der er præcis ét korrekt svar eller at der er en \emph{mængde} af ét eller flere korrekte svar.
I begge tilfælde er det tilladt at vælge mere end ét svar.
Ved retning af multiple-choice opgaver gives der negative point for forkerte svar.
Eksempler:
\begin{itemize}
\item Opgave X har 8 svarmuligheder, hvoraf 4 er korrekte. Hvert kryds ved en korrekt svarmulighed giver 1 point, og hvert kryds ved en forkert svarmulighed giver $-1$ point. Hvis du fx sætter kryds ved alle de korrekte svarmuligheder og én forkert svarmulighed får du 3 point.
\item Opgave Y har 3 svarmuligheder, hvoraf 1 er korrekt. Kryds ved den korrekte svarmulighed giver 4 point, og kryds ved en forkert svarmulighed giver $-2$ point.
\end{itemize}
Hvis du ikke sætter nogen krydser får du altid 0 point, og sætter du krydser ved alle svarmuligheder får du også altid 0 point.

\newpage

\input{\problemfile}

\end{document}
